% Part: first-order-logic
% Chapter: syntax-and-semantics
% Section: formation-sequences

\documentclass[../../../include/open-logic-section]{subfiles}

\begin{document}

\olfileid{fol}{syn}{fseq}

\olsection{متواليات التكوين}

إن تعريف ال!!{formula}s بتعريف استقرائي، والأسلوب المتمم له المتمثل
في إثبات خصائص ال!!{formula}s بالاستقراء، نهج أنيق وفعال. غير أنه قد
يكون من المفيد أيضًا النظر في نهج تصاعدي خطوة خطوة لبناء
ال!!{formula}s، وهو ما نفعله هنا باستخدام مفهوم \emph{متوالية
التكوين}.
%
ولكي نبيّن كيف يمكن إدخال الحدود وال!!{formula}s بهذه الطريقة من غير
حاجة إلى الإحالة إلى تعريفيهما الاستقرائيين، ندخل أولًا مفهوم سلسلة
اعتباطية من الرموز المأخوذة من لغة ما~$\Lang L$.

\begin{defn}[السلاسل]
\ollabel{defn:string}
افترض أن $\Lang L$ لغة رتبة أولى. \emph{سلسلة~$\Lang L$} هي
متوالية منتهية من رموز~$\Lang L$. وحيث تكون اللغة~$\Lang L$ معيّنة
بوضوح من السياق، سنشير غالبًا إلى سلسلة~$\Lang L$ باسم
\emph{سلسلة} فحسب.
\end{defn}

\begin{ex}
لكل لغة رتبة أولى~$\Lang L$، جميع
!!{formula}s-$\Lang L$ هي سلاسل~$\Lang L$، لكن العكس غير صحيح.
فمثلًا، \[)(\Obj v_0\lif\lexists\] سلسلة~$\Lang L$، لكنها ليست
!!{formula}-$\Lang L$.
\end{ex}

\begin{defn}[متواليات تكوين الحدود]
\ollabel{defn:fseq-trm}
تكون المتوالية المنتهية من سلاسل~$\Lang L$
$\tuple{t_0,\dotsc,t_n}$ \emph{متوالية تكوين} لحد~$t$ إذا كان
$t \ident t_n$، وكان، لكل $i \leq n$، إما أن $t_i$ !!a{variable}
أو !!a{constant}، وإما أن تحتوي $\Lang L$ !!{function} ذات
$k$ مواضع~$f$ وتوجد $m_1,\dotsc,m_k < i$ بحيث
$t_i \ident f(t_{m_1},\dotsc,t_{m_k})$. وحين يلزم التمييز، سنسمي
متواليات التكوين للحدود \emph{متواليات تكوين الحدود}.
\end{defn}

\begin{ex}
المتوالية
\[
    \tuple{\Obj c_0, \Obj v_0, \Atom{\Obj f^2_0}{\Obj c_0, \Obj v_0}, \Atom{\Obj f^1_0}{\Atom{\Obj f^2_0}{\Obj c_0, \Obj v_0}}}
\]
متوالية تكوين للحد~$\Atom{\Obj f^1_0}{\Atom{\Obj
f^2_0}{\Obj c_0, \Obj v_0}}$، وكذلك المتوالية
\[
    \tuple{\Obj v_0, \Obj c_0, \Atom{\Obj f^2_0}{\Obj c_0, \Obj v_0}, \Atom{\Obj f^1_0}{\Atom{\Obj f^2_0}{\Obj c_0, \Obj v_0}}}.
\]
\end{ex}

\begin{defn}[متواليات تكوين الصيغ]
\ollabel{defn:fseq-frm}
تكون المتوالية المنتهية من سلاسل~$\Lang L$
$\tuple{!A_0,\dotsc,!A_n}$ \emph{متوالية تكوين} لـ~$!A$ إذا كان
$!A \ident !A_n$، وكان، لكل $i \leq n$، إما أن $!A_i$
!!{formula} ذرية، وإما أن يوجد $j,k < i$ و!!a{variable}~$x$ بحيث
يتحقق أحد ما يأتي:
\begin{enumerate}
    \tagitem{prvNot}{$!A_i \ident \lnot !A_j$.}{}%
    \tagitem{prvAnd}{$!A_i \ident (!A_j \land !A_k)$.}{}%
    \tagitem{prvOr}{$!A_i \ident (!A_j \lor !A_k)$.}{}%
    \tagitem{prvIf}{$!A_i \ident (!A_j \lif !A_k)$.}{}%
    \tagitem{prvIff}{$!A_i \ident (!A_j \liff !A_k)$.}{}%
    \tagitem{prvAll}{$!A_i \ident \lforall[x][!A_j]$.}{}%
    \tagitem{prvEx}{$!A_i \ident \lexists[x][!A_j]$.}{}%
\end{enumerate}
وحين يلزم التمييز، سنسمي متواليات التكوين للصيغ
\emph{متواليات تكوين الصيغ}.
\end{defn}

\begin{ex}
\[
\tuple{
    \Atom{\Obj A^1_0}{\Obj v_0},
    \Atom{\Obj A^1_1}{\Obj c_1},
    (\Atom{\Obj A^1_1}{\Obj c_1} \land \Atom{\Obj A^1_0}{\Obj v_0}),
    \lexists[\Obj v_0][(\Atom{\Obj A^1_1}{\Obj c_1} \land \Atom{\Obj A^1_0}{\Obj v_0})]
}
\]
متوالية تكوين لـ~$\lexists[\Obj v_0][(\Atom{\Obj A^1_1}{\Obj
c_1} \land \Atom{\Obj A^1_0}{\Obj v_0})]$، وكذلك
\begin{multline*}
\tuple{
    \Atom{\Obj A^1_0}{\Obj v_0},
    \Atom{\Obj A^1_1}{\Obj c_1},
    (\Atom{\Obj A^1_1}{\Obj c_1} \land \Atom{\Obj A^1_0}{\Obj v_0}),
    \Atom{\Obj A^1_1}{\Obj c_1},\\
    \lforall[\Obj v_1][\Atom{\Obj A^1_0}{\Obj v_0}],
    \lexists[\Obj v_0][(\Atom{\Obj A^1_1}{\Obj c_1} \land \Atom{\Obj A^1_0}{\Obj v_0})]
}.
\end{multline*}
%
وكما يُرى من المثال الثاني، قد تحتوي متواليات التكوين على
«حشو»: !!{formula}s زائدة أو لا تسهم في البناء.
\end{ex}

\begin{prop}\ollabel{prop:formed}
لكل !!{formula}~$!A$ في~$\Frm[L]$ متوالية تكوين.
\end{prop}

\begin{proof}
افترض أن $!A$ ذرية. عندئذ تكون المتوالية $\tuple{!A}$ متوالية
تكوين لـ~$!A$.
%
وافترض الآن أن لـ~$!B$ و$!C$ متواليتي التكوين
$\tuple{!B_0,\dotsc,!B_n}$ و$\tuple{!C_0,\dotsc,!C_m}$ على
التوالي.
%
\begin{enumerate}
  \tagitem{prvNot}{إذا كان $!A \ident \lnot !B$،
      فإن $\tuple{!B_0,\dotsc,!B_n,\lnot !B_n}$
      متوالية تكوين لـ~$!A$.}{}
  \tagitem{prvAnd}{إذا كان $!A \ident (!B \land !C)$،
      فإن $\tuple{!B_0,\dotsc,!B_n,!C_0,\dotsc,!C_m,(!B_n \land !C_m)}$
      متوالية تكوين لـ~$!A$.}{}
  \tagitem{prvOr}{إذا كان $!A \ident (!B \lor !C)$،
      فإن $\tuple{!B_0,\dotsc,!B_n,!C_0,\dotsc,!C_m,(!B_n \lor !C_m)}$
      متوالية تكوين لـ~$!A$.}{}
  \tagitem{prvIf}{إذا كان $!A \ident (!B \lif !C)$،
      فإن $\tuple{!B_0,\dotsc,!B_n,!C_0,\dotsc,!C_m,(!B_n \lif !C_m)}$
      متوالية تكوين لـ~$!A$.}{}
  \tagitem{prvIff}{إذا كان $!A \ident (!B \liff !C)$،
      فإن $\tuple{!B_0,\dotsc,!B_n,!C_0,\dotsc,!C_m,(!B_n \liff !C_m)}$
      متوالية تكوين لـ~$!A$.}{}
  \tagitem{prvAll}{إذا كان $!A \ident \lforall[x][!B]$،
      فإن $\tuple{!B_0,\dotsc,!B_n,\lforall[x][!B_n]}$
      متوالية تكوين لـ~$!A$.}{}
  \tagitem{prvEx}{إذا كان $!A \ident \lexists[x][!B]$،
      فإن $\tuple{!B_0,\dotsc,!B_n,\lexists[x][!B_n]}$
      متوالية تكوين لـ~$!A$.}{}
\end{enumerate}
وبحسب مبدأ الاستقراء على ال!!{formula}s، لكل !!{formula} متوالية
تكوين.
\end{proof}

ويمكننا أيضًا إثبات العكس. وهذا مهم لأنه يبيّن أن طريقتينا في تعريف
الصيغ متكافئتان: فهما تعطيان النتائج نفسها. ويعني كذلك أن في وسعنا
إثبات مبرهنات عن الصيغ باستخدام الاستقراء العادي على طول متواليات
التكوين.

\begin{lem}
\ollabel{lem:fseq-init}
افترض أن $\tuple{!A_0,\dotsc,!A_n}$ متوالية تكوين لـ~$!A_n$، وأن
$k \leq n$. عندئذ تكون $\tuple{!A_0,\dotsc,!A_k}$ متوالية تكوين
لـ~$!A_k$.
\end{lem}

\begin{proof}
تمرين.
\end{proof}

\begin{prob}
أثبت \olref[fol][syn][fseq]{lem:fseq-init}.
\end{prob}

\begin{thm}
\ollabel{thm:fseq-frm-equiv}
$\Frm[L]$ هي مجموعة جميع سلاسل~$\Lang L$، $!A$، التي توجد لها
متوالية تكوين صيغ لـ~$!A$.
\end{thm}

\begin{proof}
لتكن $F$ مجموعة جميع سلاسل الرموز في اللغة~$\Lang L$ التي لها
متوالية تكوين. وقد رأينا في
\olref[fol][syn][fseq]{prop:formed} أن $\Frm[L] \subseteq F$،
فنثبت الآن العكس.

افترض أن لـ~$!A$ متوالية التكوين
$\tuple{!A_0,\dotsc,!A_n}$. نثبت أن $!A \in \Frm[L]$ بالاستقراء
القوي على~$n$. وفرضية استقرائنا هي أن كل سلسلة رموز لها متوالية
تكوين فهرسها الأخير $m < n$ تنتمي إلى~$\Frm[L]$. وبحسب تعريف
متوالية التكوين، إما أن تكون $!A_n$ ذرية، وإما أن يوجد $j,k < n$
و!!a{variable}~$x$ بحيث تصدق إحدى الحالات الآتية:
\begin{enumerate}
    \tagitem{prvNot}{$!A \ident \lnot !A_j$.}{}%
    \tagitem{prvAnd}{$!A \ident (!A_j \land !A_k)$.}{}%
    \tagitem{prvOr}{$!A \ident (!A_j \lor !A_k)$.}{}%
    \tagitem{prvIf}{$!A \ident (!A_j \lif !A_k)$.}{}%
    \tagitem{prvIff}{$!A \ident (!A_j \liff !A_k)$.}{}%
    \tagitem{prvAll}{$!A \ident \lforall[x][!A_j]$.}{}%
    \tagitem{prvEx}{$!A \ident \lexists[x][!A_j]$.}{}%
\end{enumerate}
نستدل الآن بالحالات. فإذا كانت $!A$ ذرية، كان
$!A_n \in \Frm[L]$. وافترض، بدلًا من ذلك، أن
$!A \ident (!A_j \land !A_k)$. بحسب
\olref[fol][syn][fseq]{lem:fseq-init}،
تكون $\tuple{!A_0,\dotsc,!A_j}$ و$\tuple{!A_0,\dotsc,!A_k}$
متواليتي تكوين لـ~$!A_j$ و$!A_k$، على التوالي. ولما كانتا متواليتين
جزئيتين ابتدائيتين حقيقيتين من متوالية تكوين~$!A$، كان الفهرس
الأخير لكل منهما أصغر من~$n$. ولذلك تنتمي $!A_j$ و$!A_k$، بفرضية
الاستقراء، إلى~$\Frm[L]$، وينتمي كذلك، بحسب تعريف !!a{formula}،
$(!A_j \land !A_k)$. وتتبع الحالات الأخرى باستدلال مماثل.
\end{proof}

لمتواليات تكوين الحدود خصائص مماثلة لخصائص متواليات تكوين
ال!!{formula}s.

\begin{prop}
\ollabel{prop:fseq-trm-equiv}
$\Trm[L]$ هي مجموعة جميع سلاسل~$\Lang L$، $t$، التي توجد لها
متوالية تكوين حدود لـ~$t$.
\end{prop}

\begin{proof}
تمرين.
\end{proof}

\begin{prob}
أثبت \olref[fol][syn][fseq]{prop:fseq-trm-equiv}.
تلميح: استخدم استراتيجية مماثلة لما استُخدم في برهان
\olref[fol][syn][fseq]{thm:fseq-frm-equiv}.
\end{prob}

ثمة نوعان من «الحشو» قد يظهران في متواليات التكوين: العناصر
المكررة، والعناصر التي لا صلة لها ببناء الصيغة أو الحد. ويمكننا
التخلص منهما معًا بالنظر في متواليات التكوين الصغرى.

\begin{defn}[متواليات التكوين الصغرى]
\ollabel{defn:minimal-fseq}
تكون متوالية التكوين $\tuple{!A_0, \ldots, !A_n}$ لصيغة~$!A$
\emph{متوالية تكوين صغرى} لـ~$!A$ إذا كان، لكل متوالية تكوين
أخرى~$s$ لـ~$!A$، طول~$s$ أكبر من أو يساوي~$n+1$.

وبالمثل، تكون متوالية التكوين $\tuple{t_0, \ldots, t_n}$
لحد~$t$ \emph{متوالية تكوين صغرى} لـ~$t$ إذا كان، لكل متوالية
تكوين أخرى~$s$ لـ~$t$، طول~$s$ أكبر من أو يساوي~$n+1$.
\end{defn}

لاحظ أنه قد يكون للصيغة أو الحد أكثر من متوالية تكوين صغرى، لكنها ستحتوي
السلاسل نفسها بالضبط.

\begin{prop}
\ollabel{prop:subformula-equivs}
الشروط الآتية متكافئة:
\begin{enumerate}
    \item $!B$ !!{subformula} من~$!A$.
    \item ترد $!B$ في كل متوالية تكوين لـ~$!A$.
    \item ترد $!B$ في متوالية تكوين صغرى لـ~$!A$.
\end{enumerate}
\end{prop}

\begin{proof}
تمرين.
\end{proof}

\begin{prob}
أثبت \olref[fol][syn][fseq]{prop:subformula-equivs}.
\end{prob}

\begin{history}
أدخل ريموند سموليان متواليات التكوين في كتابه
\emph{First-Order Logic} \citep{Smullyan1968}.
وأثبت \citet{Zuckerman1973} خصائص إضافية لمتواليات التكوين.
\end{history}

\end{document}
